%=====================================================================
\section{Fitting inside the database (v1.0)}\label{sec:v10}
Version 1.0 packages the engine for production use and fits fact tables larger than memory from statistics computed where the data lives (\emph{pushdown}). A fact table is \emph{pushed down} when it is a fact that no table references, has no type~2 history, and is not the source of an aggregate column that the source does not store (\texttt{relsynth/pushdown/planner.py}). Its rows never leave the database as a whole: the engine reads (i) exact aggregates computed by \texttt{GROUP BY} and window functions, and (ii) a uniform sample of at most \texttt{synthesis.pushdown.sample\_rows} rows. Every other table is loaded as before. The propositions below show that this changes nothing in the relational model, and that the attribute models fitted on the sample are an exact posterior given the sampled rows. The checks are in \texttt{tests/unit/test\_sqlfuncs\_stats.py}, \texttt{tests/unit/test\_pushdown.py} and \texttt{tests/integration/test\_servers.py}; they run on SQLite, PostgreSQL~16, SQL Server~2022 and Oracle Database~26ai Free (23.26).

\subsection{What the relational model needs from a fact}
Let a pushed fact $F$ have $N$ rows, key columns $k_1..k_m$ (links to dimensions or to a parent fact, the parts of a composite key) and continuous columns $c_1..c_p$. For a non-null value $x_i$ of column $c$ with average rank $r_i$ among the $N_c$ non-null values let $z_i=\Phi^{-1}\bigl(r_i/(N_c+1)\bigr)$ (its \emph{normal score}), and $z_i=0$ for a null. For a key tuple $u$ write
\[n(u)=\#\{i:k(i)=u\},\qquad S_c(u)=\sum_{i:k(i)=u}z_i^{(c)},\qquad Q_c(u)=\sum_{i:k(i)=u}\bigl(z_i^{(c)}\bigr)^2,\]
with a null key value forming a group of its own.

\begin{proposition}[Sufficiency of the in-database statistics]\label{prop:pdsuff}
Given the loaded tables (entities and non-pushed facts), the full conditional of every latent quantity of the relational model of Sections~\ref{sec:star}--\ref{sec:v09} that involves a pushed fact $F$---segment assignments of the entities it links, block and link rates (Props.~\ref{prop:eta}, \ref{prop:split}), count and multiplicity laws (Props.~\ref{prop:mult}, \ref{prop:countbf}, \ref{prop:dpclosed}), participation and inclusion (Props.~\ref{prop:mandatory}, \ref{prop:particippool}, \ref{prop:sparsedir}), nested refinement (Prop.~\ref{prop:nested}), order headers (Prop.~\ref{prop:header}) and running balances (Prop.~\ref{prop:cumulative})---depends on the rows of $F$ only through
\begin{enumerate}
\item the counts $n(u)$ over the key tuples the model uses (pairs of links, a composite key with its links, a header key with its header columns);
\item the score sums $S_c(u)$ and $Q_c(u)$ for the (at most three) continuous columns that segment discovery uses;
\item for an order header, the first row of every document key in primary-key order and the number of cells that differ from it; for a running balance, the first row of every group in the rule's order, the number of non-null values and the minimum;
\item $N$, the number of non-null values of each key, and the distinct referenced values of each key.
\end{enumerate}
\end{proposition}
\begin{proof}
Each listed model block is a product over the rows of $F$ of a factor that depends on a row only through its key tuple $u$ and, for segment discovery, through the rank features $z^{(c)}$ of the row. Block-rate and multiplicity likelihoods are products of Poisson or multinomial terms whose exponents are the counts of rows per pair of linked rows, so they are functions of $n(u)$ (the conjugate posteriors of Props.~\ref{prop:eta} and \ref{prop:mult} use sums of $n(u)$ over segment pairs). The rank-feature factor of segment discovery is Gaussian in the scores, with a mean per pair of segments and a common variance per column (\texttt{PairRel}, \texttt{OccRel}), so its log-likelihood over the rows of a key tuple is $a\,n(u)+b\,S_c(u)+c\,Q_c(u)$ for coefficients that do not depend on the rows, and the conjugate updates of the means and variances use the same three sums. Participation and inclusion use which keys occur (the distinct values) and how often ($n(u)$). A nested fact is linked to its parent fact through its composite key, so the counts over (composite key, links) give the parent row and the dimension rows of every group, with the multiplicity $n(u)$ as the group's weight. An order header is the first row of its document key (Prop.~\ref{prop:header}); its split is lossless exactly when no cell differs from that row, which is what (3) counts. A running balance needs its opening value per group (the first row in the rule's order) and its floor (the minimum); the rest of the column is a deterministic function of the flows (Prop.~\ref{prop:cumulative}). Nothing else of $F$ enters these factors.
\end{proof}

\begin{proposition}[Agreement with in-memory fitting]\label{prop:pdexact}
Let the loaded rows of every table be put in a canonical order: by primary key, else by the first declared unique key that holds, and by a hash of the row's values when there is neither (\texttt{canonical\_order}); and let order headers split from a fact be put in the order of their document key. If the sample of a pushed fact holds all its rows, the fitted model and every synthetic dataset are identical, value for value, to those of fitting the whole table in memory with the same seed; and the model fitted from a database does not depend on the order in which the database returns rows.
\end{proposition}
\begin{proof}
By Prop.~\ref{prop:pdsuff} the relational blocks see the same numbers, computed either by \texttt{FrameStats} (pandas) or by \texttt{SqlStats} (SQL); the attribute models see the same rows. The fit is a deterministic function of its inputs, the seed and the order in which rows are visited (the Gibbs sweeps visit rows in order). The canonical order (by primary key, else by the first declared unique key, else by a hash of the row; order headers split from a fact by their document key) is a function of the set of rows only, so two backends, or two scans of one table, give the same input sequence. The only remaining difference is floating-point: SQL sums are accumulated in another order than in pandas, and servers compute $\Phi^{-1}$ with Prop.~\ref{prop:acklam}; both perturb $S_c,Q_c$ by relative errors below $10^{-8}$, far under the resolution at which a segment draw changes. The checks compare the synthetic tables cell by cell: identical for all four pushed fact shapes on SQLite (\texttt{test\_full\_sample\_pushdown\_equals\_in\_memory}) and for a pushed star schema on PostgreSQL, SQL Server and Oracle against SQLite (\texttt{test\_full\_sample\_pushdown\_matches\_sqlite\_in\_memory}).
\end{proof}
\noindent Without the canonical order the same database gave different models on different servers (SQL Server and Oracle return rows in a different order than SQLite and PostgreSQL), and a server using parallel scans could give different models on two runs. v1.0 orders every loaded table; this changes in-memory results relative to v0.9 by a permutation of the input rows, which leaves the posterior unchanged and changes individual draws only.

\subsection{Normal scores on database servers}
\begin{proposition}[Acklam's quantile approximation]\label{prop:acklam}
PostgreSQL, SQL Server and Oracle have no standard-normal quantile function. relsynth computes $\Phi^{-1}(p)$ in SQL with Acklam's rational approximation (a rational function of $q=p-\tfrac12$ on $[0.02425, 0.97575]$ and of $\sqrt{-2\log p}$ in the tails), whose relative error is below $\delta=1.15\times10^{-9}$ on $(0,1)$. Hence $|\hat S_c(u)-S_c(u)|\le \delta\sum_{i:k(i)=u}|z_i|$ and $|\hat Q_c(u)-Q_c(u)|\le (2\delta+\delta^2)\,Q_c(u)$. SQLite registers the exact function (\texttt{scipy.special.ndtri}) as a user function.
\end{proposition}
\begin{proof}
The error bound of the approximation is Acklam's (2003), and \texttt{test\_acklam\_error\_bound} confirms it in double precision on 4{,}081 points spanning $[10^{-12},1-10^{-12}]$ (relative error where $|\Phi^{-1}(p)|>10^{-3}$, absolute error below $10^{-8}$ everywhere). For a sum, $|\sum\hat z_i-\sum z_i|\le\sum|\hat z_i-z_i|\le\delta\sum|z_i|$ with $\delta=1.15\times10^{-9}$; for squares $|\hat z^2-z^2|=|z|^2\,|(1+\epsilon)^2-1|\le(2\delta+\delta^2)z^2$.
\end{proof}

\subsection{The sample}
\begin{proposition}[Seeded hash sampling is a Bernoulli sample]\label{prop:hashsample}
The sample keeps row $i$ when $h(s,\kappa_i)<\lfloor q\,2^{32}\rfloor$, where $\kappa_i$ is the row's primary key (all its columns when there is none), $s$ the seed, $q=\min(1,n/N)$ for the requested size $n$, and $h$ the database's hash reduced to $[0,2^{32})$ (\texttt{hashtextextended} on PostgreSQL, SHA-256 by \texttt{HASHBYTES} on SQL Server, \texttt{ORA\_HASH} on Oracle, BLAKE2b on SQLite and files). If $h(s,\cdot)$ behaves as a uniform random function independent of the data, then
(a) the inclusion indicators are i.i.d.\ Bernoulli$(q)$ and independent of every column value, so the sample is a uniform random subset with $|S|\sim\mathrm{Bin}(N,q)$;
(b) the same seed gives the same sample on every run and every engine with the same hash;
(c) $P\bigl(\bigl||S|-Nq\bigr|>8\sqrt{Nq(1-q)}+10\bigr)<1.3\times10^{-15}$, so the guard that stops a sample of that size (a broken hash, such as one returning NULL) does not stop a valid sample in practice.
\end{proposition}
\begin{proof}
(a) Under the assumption, $h(s,\kappa_i)$ are i.i.d.\ uniform on $[0,2^{32})$ because the keys are distinct, so each comparison is Bernoulli with probability $\lfloor q2^{32}\rfloor/2^{32}$, which differs from $q$ by less than $2^{-32}$. (b) The comparison is a deterministic function of $(s,\kappa_i)$. (c) With $\sigma^2=Nq(1-q)$ and $t=8\sigma+10$, Bernstein's inequality gives $P(||S|-Nq|>t)\le 2\exp\bigl(-t^2/(2\sigma^2+2t/3)\bigr)\to2e^{-32}<2.6\times10^{-14}$ as $\sigma\to\infty$; the exact binomial tails, computed for $N\in\{10,\dots,10^9\}$ and $q\in[10^{-6},1-10^{-6}]$, are at most $1.24\times10^{-15}$, the normal value $2\Phi(-8)$ approached for large $\sigma$ (for small $\sigma$ the additive 10 makes the tail far smaller). \texttt{test\_python\_hash\_is\_uniform\_and\_seeded} tests uniformity, and \texttt{test\_sample\_is\_uniform\_reproducible\_and\_sized} size and reproducibility on each engine.
\end{proof}

\begin{proposition}[Attribute models from the sample are an exact posterior]\label{prop:pdsample}
Let the sample be drawn as in Prop.~\ref{prop:hashsample}. The posterior of the attribute-model parameters $\theta$ of the pushed fact given the sampled rows,
\[p(\theta\mid x_S,S)\propto p(\theta)\prod_{i\in S}p(x_i\mid\theta,\text{keys}_i),\]
is the exact posterior given the observed data: the sampling mechanism is ignorable. Synthetic datasets are therefore posterior-predictive draws (Prop.~\ref{prop:postpred}) given the sampled rows for the attributes and given all rows for the relational part; their sizes follow the exact $N$, not $|S|$. The posterior is wider than with all rows (its variance scales with $1/|S|$ for regular parameters), so the scorecard's real side uses exact statistics where they exist (row counts, null keys, cardinalities, duplicate rates, pair multiplicities) and the sample elsewhere, with limits computed for the sample's size.
\end{proposition}
\begin{proof}
The inclusion indicators are independent of the values and of $\theta$ (Prop.~\ref{prop:hashsample}a), so $p(S\mid x,\theta)=p(S)$ and the likelihood of the observed data factors as $p(S)\prod_{i\in S}p(x_i\mid\theta)$; the factor $p(S)$ does not involve $\theta$ (Rubin 1976: data missing completely at random). Prop.~\ref{prop:pdsuff} makes the relational blocks depend on all rows through exact statistics. The sizes use $N$ because the scale rule (Prop.~\ref{prop:scale}) is applied to the exact counts.
\end{proof}

\subsection{Why entities are always loaded}
\begin{proposition}[Eligibility]\label{prop:pdeligible}
A table cannot be fitted from aggregates alone when any of the following holds: it is an entity (each of its rows carries its own latent segment, whose full conditional needs that row's attributes and links, Prop.~\ref{prop:perrow}); another table references it (it is then an entity of its own rows for that table's links, including the parent of a nested fact); it has type~2 history (Prop.~\ref{prop:scd2} draws versions per business key and row); or an aggregate column of another table is computed from its rows and the source does not store that column (the column must be recomputed from the rows). Every other fact is eligible, including nested facts, facts split into order headers, bases of views and facts with running balances.
\end{proposition}
\begin{proof}
For each excluded case the named proposition's full conditional has a factor per row of the table that is not a function of the statistics of Prop.~\ref{prop:pdsuff}. For the eligible shapes, Prop.~\ref{prop:pdsuff} lists the statistics that replace the rows.
\end{proof}

\noindent\emph{Memory and time.} A pushed fact costs memory proportional to the number of distinct key tuples plus the sample, instead of rows times columns: on AdventureWorks (all ten eligible facts pushed, including the 776{,}286-row inventory fact with its running balance) the peak resident memory fell from 2.1\,GiB to 1.6\,GiB at a sample of 30{,}000 rows, and the run took 668\,s instead of 803\,s on SQLite. The database does the grouping; the number of round trips is a constant per statistic and does not grow with $N$.

\subsection{Sizes of child entities}
An entity $C$ whose rows link to a parent entity $P$ (orders to customers) has its own size factor $s_C$; the link model draws each of its rows' parents from the segments of $P$. If only $P$ is resized, the expected number of $C$ rows per $P$ row changes by $s_C/s_P$, which is rarely what a user who asks for ``ten times the customers'' wants. From v1.0 an entity without a factor of its own takes the factor of the first parent it links to that has one (\texttt{inherit\_scale}, applied in topological order, so a chain of entities follows its root), and the scorecard expects the same sizes. With $s_C=s_P$ the expected rows per parent are those of the source, and facts anchored on $C$ follow by Prop.~\ref{prop:scale}. A table given its own factor keeps it. \emph{Test} \texttt{test\_child\_entities\_follow\_a\_scaled\_parent}: customers $\times3$ gives exactly $3\times$ the invoices and $3\times$ the lines within 10\%, with every invoice-per-customer check passing.
